\documentclass[12pt]{article}
\usepackage[utf8]{inputenc}
\usepackage[margin=1in]{geometry}
\usepackage{amsmath}
\usepackage{graphicx}
\usepackage{booktabs}
\usepackage{hyperref}
\usepackage[round]{natbib}

\title{Glutamate Neurotransmission from Leptin Receptor Cells Is Required for Typical Puberty and Reproductive Function in Female Mice}
\author{Author A$^{1}$, Author B$^{1}$, Author C$^{2}$, Author D$^{1,2}$ \\
\small $^{1}$Department of Molecular and Integrative Physiology, University of Michigan, Ann Arbor, MI \\
\small $^{2}$Neuroscience Graduate Program, University of Michigan, Ann Arbor, MI}
\date{}

\begin{document}
\maketitle

\begin{abstract}
Leptin is an essential metabolic signal for the reproductive axis, yet the neural circuit mechanisms through which it controls puberty onset and fertility remain incompletely understood. The ventral premammillary nucleus (PMv) is a glutamatergic hypothalamic nucleus densely expressing the leptin receptor (LepRb) that projects directly to gonadotropin-releasing hormone (GnRH) neurons, but its functional contribution has been overshadowed by studies emphasizing GABAergic pathways. Here, we used three complementary approaches in female mice to demonstrate that glutamatergic neurotransmission from LepRb-expressing PMv neurons is required for leptin's control of reproduction. Chemogenetic activation of PMv LepRb neurons induced transient luteinizing hormone (LH) release in fed adult females, with LH levels correlating with the number of activated PMv neurons ($n = 15$ PMv-hit animals). Conditional deletion of Vglut2 in LepRb neurons (LepRb-$\Delta$Vglut2) delayed the age of first estrus ($p = 0.028$), disrupted estrous cyclicity, and increased GnRH immunoreactivity in median eminence terminals ($p = 0.0005$), suggesting impaired GnRH release. In a PMv-specific rescue model, restoring LepRb expression in the PMv of otherwise leptin-receptor-null mice rescued vaginal opening, follicle maturation, and pregnancy, whereas simultaneous deletion of Vglut2 in the PMv completely prevented pubertal development. These findings establish that glutamate release from LepRb PMv neurons is necessary for leptin's regulation of the reproductive axis, challenging the prevailing view that only GABAergic leptin signaling is relevant for reproduction.
\end{abstract}

\section{Introduction}

The adipocyte-derived hormone leptin serves as a critical metabolic signal that permits activation of the hypothalamic-pituitary-gonadal axis \citep{elias2012leptin}. In both humans and rodents, leptin deficiency results in infertility and delayed puberty, underscoring the importance of adequate energy stores for reproductive competence \citep{clarke2015effects}. Despite decades of research, the precise neural circuits and neurotransmitter mechanisms through which leptin controls gonadotropin-releasing hormone (GnRH) secretion and downstream reproductive function have remained a subject of active investigation \citep{herbison2016gnrh}.

The leptin receptor (LepRb) is expressed in multiple hypothalamic populations, including neurons in the arcuate nucleus, the ventromedial hypothalamus, and the ventral premammillary nucleus (PMv) \citep{leshan2009leptin}. Among these, the PMv has emerged as a particularly compelling candidate for mediating leptin's reproductive effects because it sends dense, direct projections to GnRH neuronal cell bodies in the preoptic area and to kisspeptin neurons in the anteroventral periventricular nucleus (AVPV) and arcuate nucleus \citep{donato2013pmv}. The PMv is a predominantly glutamatergic nucleus, with the majority of its LepRb-positive neurons co-expressing the vesicular glutamate transporter Vglut2 (Slc17a6) \citep{donato2011ventral}.

Previous studies using conditional deletion of the leptin receptor from all Vglut2-expressing neurons showed surprisingly mild reproductive phenotypes, whereas deletion of LepRb from GABAergic neurons (using Vgat-Cre) produced profound reproductive deficits \citep{zuure2013leptin, vong2011leptin}. These findings led to the prevailing conclusion that GABAergic, rather than glutamatergic, neurotransmission mediates the critical reproductive actions of leptin. However, the global Vglut2-neuron deletion strategy is inherently broad, affecting LepRb signaling in numerous glutamatergic populations simultaneously, and may therefore mask the specific contribution of PMv glutamate signaling through compensatory mechanisms in other cell populations.

To address this gap, we employed three complementary strategies. First, we used chemogenetic activation (Designer Receptors Exclusively Activated by Designer Drugs, DREADDs) to determine whether acute stimulation of PMv LepRb neurons can induce LH secretion in fed adult female mice \citep{roth2016dreadds}. Second, we generated conditional knockout mice lacking Vglut2 specifically in LepRb neurons (LepRb-$\Delta$Vglut2) to assess the consequences for pubertal development, estrous cyclicity, GnRH content, and LH dynamics. Third, we used a PMv-specific rescue model (LepRloxTB) in which LepRb expression is restored exclusively in the PMv by stereotaxic delivery of AAV-Cre, with or without concomitant deletion of Vglut2, to determine whether glutamate release from PMv neurons is specifically required for reproductive rescue. Our results demonstrate that glutamatergic neurotransmission from LepRb-expressing PMv neurons is essential for normal puberty and reproductive function, establishing the PMv as a key node in leptin's control of the reproductive axis.

\section{Methods}

\subsection{Animals}

All procedures were approved by the University of Michigan Institutional Animal Care and Use Committee (protocol PRO00010420). Mice were housed under a 12:12 light-dark cycle (lights on at 06:00) at 21--23$^\circ$C with ad libitum access to food and water. Maintenance diet was Envigo 2016 (low phytoestrogen); breeding diet was Envigo 2019 Teklad. The transgenic lines used are summarized in Table~\ref{tab:models}.

\begin{table}[htbp]
\centering
\caption{Mouse models used in this study.}
\label{tab:models}
\begin{tabular}{llll}
\toprule
Model & Genotype & JAX Stock & Purpose \\
\midrule
LepRb-Cre & LeprCre knock-in & 032457 & Cre in LepRb neurons \\
LepRloxTB & Lepr-null, Cre-rescue & 018989 & PMv-specific LepRb restoration \\
Vglut2-floxed & Slc17a6 floxed & 012898 & Conditional Vglut2 deletion \\
Kiss1-hrGFP & Kiss1-hrGFP reporter & 023425 & Kisspeptin neuron visualization \\
C57BL/6 & Wild-type & 000664 & Controls \\
\bottomrule
\end{tabular}
\end{table}

\subsection{Viral Vectors and Stereotaxic Surgery}

Four viral constructs were used: AAV8-hSyn-DIO-hM3Dq-mCherry (Addgene \#44361) for stimulatory DREADD expression, AAV8-hSyn-DIO-mCherry (Addgene \#50459) as a control reporter, AAV8-CMV-HI-eGFP-Cre (UPenn Vector Core) for Cre-mediated recombination, and AAV2/7-CMV-PI-EGFP (UPenn Vector Core) as a GFP control. Stereotaxic injections targeted the PMv at coordinates AP $-5.4$ mm from the rostral rhinal vein, ML $\pm 0.5$ mm from the superior sagittal sinus, and DV $-5.4$ mm from the dura mater. Approximately 50 nL of virus was delivered per side via a glass micropipette and picospritzer. Animals recovered for one month before experimentation.

\subsection{Chemogenetic Activation and Serial Blood Sampling}

LepRb-Cre mice received unilateral injections of AAV-hM3Dq into the PMv. After four weeks, animals underwent cannulation of the aortic artery (for blood sampling) and jugular vein (for intravenous drug delivery). Following three days of recovery, serial blood sampling was performed using the Culex Automated Blood Sampling System. Baseline LH was measured from three samples collected every 10 minutes. Clozapine-N-oxide (CNO, 0.5 mg/kg) or clozapine (0.5 mg/kg) was then administered intravenously in 50 $\mu$L, and LH was sampled every 10 minutes for one hour. PMv targeting was verified post hoc using mCherry and cFos immunoreactivity, and animals were classified as PMv-hits or PMv-misses accordingly.

\subsection{Conditional Knockout: LepRb-$\Delta$Vglut2}

LepRb-$\Delta$Vglut2 mice were generated by crossing LepRb-Cre with Vglut2-floxed mice. Reproductive phenotyping included monitoring the age of vaginal opening, age of first estrus, estrous cyclicity via vaginal cytology, basal LH levels, LH surge assessment, and quantification of GnRH immunoreactivity in the arcuate nucleus and median eminence using fluorescence microscopy.

\subsection{PMv-Specific Rescue Model}

LepRloxTB mice (leptin receptor null with Cre-inducible restoration) were crossed with Vglut2-floxed mice. Stereotaxic injection of AAV-Cre into the PMv simultaneously restored LepRb expression and deleted Vglut2 locally. Control animals received AAV-Cre without carrying the Vglut2-floxed allele. Pubertal development (vaginal opening), ovarian histology, and fertility (breeding tests) were assessed. Leptin responsiveness was confirmed by pSTAT3 immunoreactivity, and Vglut2 deletion was verified by in situ hybridization.

\subsection{Statistical Analysis}

Data were analyzed using Student's t-tests for two-group comparisons and Kruskal-Wallis tests with appropriate post hoc corrections for multi-group comparisons. Statistical significance was set at $\alpha = 0.05$. All values are reported as mean $\pm$ SEM unless otherwise indicated.

\section{Results}

\subsection{Chemogenetic Activation of PMv LepRb Neurons Induces LH Release}

To determine whether acute activation of LepRb neurons in the PMv can stimulate the reproductive axis, we expressed the stimulatory DREADD hM3Dq in LepRb-Cre neurons via unilateral PMv injection and monitored LH secretion following intravenous CNO administration. Of 18 total animals receiving CNO, 15 showed successful PMv targeting (PMv-hits) based on mCherry expression and cFos immunoreactivity, 2 were classified as PMv-misses (viral expression outside the PMv), and 1 was excluded due to minimal viral expression (Table~\ref{tab:dreadd}).

\begin{table}[htbp]
\centering
\caption{Experimental groups for chemogenetic activation experiments.}
\label{tab:dreadd}
\begin{tabular}{lllll}
\toprule
Group & AAV & Drug (i.v.) & $N$ & LH Response \\
\midrule
PMv-hits (CNO) & hM3Dq & CNO (0.5 mg/kg) & 15 & Variable \\
PMv-misses (CNO) & hM3Dq & CNO & 2 & No increase \\
mCherry control & mCherry & CNO & 5 & No increase \\
PMv-hits (clozapine) & hM3Dq & Clozapine (0.5 mg/kg) & 4 & Increase \\
Clozapine alone & None & Clozapine & 7 & No increase \\
\bottomrule
\end{tabular}
\end{table}

Among PMv-hit animals, a subset showed transient LH elevations following CNO, with peak values ranging from 0.65 to 3.11 ng/mL. Critically, LH levels correlated with the number of cFos-immunoreactive neurons in the PMv, indicating that the magnitude of neuronal activation predicted the endocrine response. Animals classified as PMv-misses and mCherry controls showed no LH increase, validating the spatial specificity of the effect. Furthermore, clozapine (a CNO metabolite that can directly activate DREADDs) also induced LH release in all four hM3Dq-expressing animals tested, whereas clozapine alone in animals without DREADD expression had no effect, confirming DREADD-mediated activation.

An additional observation was that some PMv-hit animals showed viral spread to the posterior arcuate nucleus (Arc), which contains POMC neurons. Analysis of cFos/mCherry colocalization in the posterior Arc revealed a significant difference among groups (Kruskal-Wallis $H = 14.63$, $p < 0.0001$), with PMv-hit animals showing LH increase having the highest Arc activation. This suggests that POMC neuron activation may mildly contribute to LH release, although the primary effect was driven by PMv targeting.

\subsection{Validation of Vglut2 Deletion in LepRb Neurons}

In situ hybridization for Lepr and Vglut2 mRNA confirmed successful deletion in LepRb-$\Delta$Vglut2 mice. In control animals, the ventromedial hypothalamus (VMHdm) and PMv showed robust colocalization of Lepr and Vglut2 transcripts. In LepRb-$\Delta$Vglut2 mice, Vglut2 expression was abolished in Lepr-positive neurons in both regions. Adult LepRb-$\Delta$Vglut2 animals displayed an obese phenotype, consistent with the known role of glutamatergic leptin signaling in energy balance regulation.

\subsection{Delayed Puberty and Disrupted Cyclicity in LepRb-$\Delta$Vglut2 Females}

We assessed pubertal milestones in LepRb-$\Delta$Vglut2 females compared with controls. While the age of vaginal opening was not significantly different between genotypes ($t_{18} = 0.84$, $p = 0.41$), the age of first estrus was significantly delayed in LepRb-$\Delta$Vglut2 females ($t_{14} = 2.45$, $p = 0.028$; Table~\ref{tab:puberty}). Body mass at both milestones did not differ significantly. In males, the age of balanopreputial separation showed a trend toward delay ($t_{11} = 1.80$, $p = 0.10$), while body mass at this milestone was significantly increased ($t_{11} = 2.43$, $p = 0.033$).

\begin{table}[htbp]
\centering
\caption{Pubertal development and statistical comparisons in LepRb-$\Delta$Vglut2 mice.}
\label{tab:puberty}
\begin{tabular}{lcccc}
\toprule
Metric & $t$-statistic & $df$ & $p$-value & Significance \\
\midrule
Age of vaginal opening & 0.84 & 18 & 0.41 & NS \\
Body mass at VO & 1.47 & 18 & 0.16 & NS \\
Age of first estrus & 2.45 & 14 & 0.028 & * \\
Body mass at first estrus & 0.59 & 14 & 0.56 & NS \\
Age of BPS (males) & 1.80 & 11 & 0.10 & NS \\
Body mass at BPS (males) & 2.43 & 11 & 0.033 & * \\
GnRH-ir in Arc/ME & 8.01 & 5 & 0.0005 & *** \\
\bottomrule
\end{tabular}
\end{table}

Estrous cyclicity monitoring revealed that control females exhibited regular 4--5 day cycles, whereas LepRb-$\Delta$Vglut2 females displayed irregular cycles with prolonged periods in diestrus. This disruption in cyclicity is consistent with impaired hypothalamic drive to the reproductive axis.

\subsection{Increased GnRH Content in Median Eminence Terminals}

Fluorescence microscopy quantification of GnRH immunoreactivity in the arcuate nucleus and median eminence revealed a striking increase in LepRb-$\Delta$Vglut2 females compared with controls ($t_5 = 8.01$, $p = 0.0005$; $n = 4$ controls, $n = 3$ mutants). This accumulation of GnRH in axon terminals is interpreted as reflecting impaired GnRH secretion rather than increased synthesis, as reduced release would lead to peptide buildup in terminal boutons. Consistent with this interpretation, LepRb-$\Delta$Vglut2 females showed altered basal LH levels during diestrus and impaired LH surge capacity, indicating a functional deficit in GnRH release.

\subsection{PMv-Specific Rescue Requires Glutamate for Reproductive Competence}

The most definitive evidence for the requirement of glutamatergic signaling came from the PMv-specific rescue experiments using the LepRloxTB model. In this system, LepRb expression is absent throughout the brain until Cre recombinase is delivered, enabling site-specific restoration. Injection of AAV-Cre into the PMv of LepRloxTB mice restored leptin signaling as confirmed by pSTAT3 immunoreactivity. These LepRb-rescued animals showed vaginal opening, follicle maturation with corpora lutea, restored estrous cycles, and became pregnant when paired with males.

In contrast, when AAV-Cre was injected into the PMv of LepRloxTB;Vglut2$^{flox/flox}$ mice---simultaneously restoring LepRb and deleting Vglut2---the animals showed restored pSTAT3 signaling but decreased Vglut2 mRNA in the PMv. Critically, these animals failed to undergo vaginal opening, showed no follicle maturation (absence of mature follicles and corpora lutea), did not develop estrous cycles, and could not become pregnant (Table~\ref{tab:rescue}).

\begin{table}[htbp]
\centering
\caption{Reproductive outcomes in PMv-specific rescue experiments.}
\label{tab:rescue}
\begin{tabular}{lcccc}
\toprule
Genotype + AAV-Cre & Vaginal Opening & Estrous Cycles & Mature Follicles & Pregnancy \\
\midrule
LepRloxTB (rescue) & Yes & Restored & Present & Yes \\
LepRloxTB;Vglut2$^{fl}$ (rescue + KO) & No & None & Absent & No \\
\bottomrule
\end{tabular}
\end{table}

The completeness of the reproductive failure in the dual rescue/knockout model provides compelling evidence that glutamate release from PMv neurons is not merely modulatory but is required for leptin-driven pubertal initiation and fertility.

\section{Discussion}

The present study provides three convergent lines of evidence that glutamatergic neurotransmission from LepRb-expressing neurons in the ventral premammillary nucleus is essential for normal puberty and reproductive function in female mice. Our chemogenetic activation experiments demonstrate that acute stimulation of PMv LepRb neurons is sufficient to induce transient LH secretion in fed adult females, with the magnitude of the response correlating with the extent of neuronal activation. The conditional deletion of Vglut2 from LepRb neurons delays pubertal maturation, disrupts estrous cyclicity, and impairs GnRH release as evidenced by peptide accumulation in median eminence terminals. Most compellingly, PMv-specific restoration of LepRb in otherwise receptor-null mice fully rescues reproductive function, but this rescue is completely abolished when Vglut2 is simultaneously deleted in the PMv.

\subsection{Reconciling with Previous Studies}

Our findings challenge the prevailing interpretation that GABAergic, rather than glutamatergic, leptin signaling is the primary pathway for reproductive control. The influential study by \citet{zuure2013leptin} demonstrated that deletion of LepRb from all Vglut2-expressing neurons produced minimal reproductive deficits, while deletion from GABAergic neurons caused severe impairment. Similarly, \citet{vong2011leptin} reported that leptin action on GABAergic neurons prevents obesity and modulates POMC neuron activity. However, the global Vglut2-neuron deletion strategy affects LepRb signaling in numerous glutamatergic populations simultaneously, including neurons in the ventromedial hypothalamus, lateral hypothalamus, and other regions. Compensatory mechanisms among these diverse populations may mask the specific contribution of PMv glutamate signaling to reproduction.

Our PMv-specific approach circumvents this limitation. By restricting both the rescue and the deletion to the PMv using stereotaxic viral delivery, we isolated the contribution of this single nucleus. The complete failure of reproductive development when Vglut2 is deleted specifically in the PMv, despite intact leptin signaling (as confirmed by pSTAT3), demonstrates that the glutamatergic output of PMv neurons is indispensable for leptin's reproductive actions.

\subsection{PMv as a Critical Relay in the Leptin-Reproduction Circuit}

The PMv occupies a strategic position in the neural circuitry linking metabolic status to reproductive competence \citep{donato2011ventral, donato2013pmv}. It receives dense leptin receptor-mediated input and sends direct projections to multiple targets critical for GnRH secretion, including GnRH cell bodies in the preoptic area, kisspeptin neurons in the AVPV, and kisspeptin/neurokinin B/dynorphin (KNDy) neurons in the arcuate nucleus \citep{cravo2011characterization, padilla2017kisspeptin}. Our observation that chemogenetic activation of PMv LepRb neurons induces LH release in fed animals, when leptin levels are presumably adequate, suggests that the PMv can serve as a positive drive signal to the GnRH system independent of acute metabolic deficit.

The finding that LH response magnitude correlated with the number of cFos-positive neurons in the PMv supports a graded, population-level encoding of the metabolic-to-reproductive signal. Interestingly, some animals with confirmed PMv targeting did not show LH elevation, potentially due to variable viral transduction efficiency, differences in CNO brain penetration, or individual variation in PMv circuit connectivity. The observation that clozapine, a known CNO metabolite, also effectively induced LH release in DREADD-expressing animals confirms that the pharmacological activation was mediated through the expressed receptors.

\subsection{Impaired GnRH Release Rather Than Synthesis}

The striking increase in GnRH immunoreactivity in median eminence terminals of LepRb-$\Delta$Vglut2 females ($p = 0.0005$) is most parsimoniously explained by impaired GnRH secretion rather than increased GnRH synthesis \citep{herbison2016gnrh, manfredi2019gnrh}. When GnRH release is compromised, the peptide accumulates in axon terminals that continue to receive newly synthesized peptide from the cell body via axonal transport. This accumulation, combined with disrupted estrous cyclicity and altered LH dynamics, paints a consistent picture of a hypothalamic deficit in the final common pathway driving pulsatile and surge-mode GnRH secretion.

\subsection{Potential Contribution of Arcuate Nucleus Neurons}

Our analysis of viral spread revealed that some animals had AAV expression extending to the posterior arcuate nucleus, which contains POMC neurons that can influence the reproductive axis. The significant difference in cFos/mCherry colocalization in the posterior Arc among experimental groups suggests that POMC neuron activation may contribute to the LH response observed in some PMv-hit animals. Importantly, KNDy neurons in the arcuate nucleus were not activated by this chemogenetic paradigm, likely because they represent a different neuronal population from the POMC neurons in this region. Future experiments using more spatially restricted viral approaches or intersectional strategies will be needed to fully parse the relative contributions of PMv versus Arc neurons.

\subsection{Implications for Understanding Metabolic Control of Fertility}

Our results have broader implications for understanding how the brain integrates metabolic information with reproductive output. The identification of glutamate as a required neurotransmitter from PMv neurons suggests that fast excitatory synaptic transmission, rather than solely neuromodulatory mechanisms, mediates the metabolic gating of puberty and cyclicity. This is consistent with the need for rapid, temporally precise signaling to coordinate pulsatile GnRH release with appropriate metabolic conditions \citep{elias2012leptin, clarke2015effects}.

The obese phenotype of LepRb-$\Delta$Vglut2 mice confirms that the same glutamatergic leptin-responsive neurons also participate in energy balance regulation, consistent with the known dual role of the PMv and VMH in both metabolic and reproductive control. This anatomical overlap provides a potential mechanism for the tight coupling between body weight regulation and reproductive competence observed across mammalian species.

\subsection{Limitations}

Several limitations should be noted. First, the LepRb-$\Delta$Vglut2 model deletes Vglut2 from all LepRb-expressing glutamatergic neurons, not exclusively those in the PMv. While the PMv-specific rescue experiments address this concern, the conditional knockout results should be interpreted as reflecting the combined loss of glutamate signaling across multiple LepRb populations. Second, our chemogenetic approach used unilateral PMv injections, and bilateral activation might produce more robust or different LH responses. Third, the LepRloxTB rescue model involves developmental leptin receptor deficiency followed by postnatal restoration, which may not perfectly recapitulate normal development. Fourth, we focused primarily on female reproductive function; the male phenotype warrants further investigation.

\section{Conclusion}

This study establishes that glutamatergic neurotransmission from leptin receptor-expressing neurons in the ventral premammillary nucleus is required for normal pubertal development and reproductive function in female mice. Acute activation of PMv LepRb neurons induces LH secretion in a dose-dependent manner correlating with neuronal activation. Loss of glutamate signaling from LepRb neurons delays puberty, disrupts cyclicity, and impairs GnRH release. PMv-specific rescue of leptin signaling fully restores reproductive competence, but only when glutamatergic transmission is intact. These findings redefine our understanding of the neurotransmitter basis of leptin's reproductive actions and identify the PMv glutamatergic output as an essential component of the metabolic-reproductive relay circuit.

\bibliographystyle{plainnat}
\bibliography{references}

\end{document}
